\section{结论}
\label{chapter:05}

% block-id: 05.conclusions.p1
\hypertarget{05.conclusions.p1}{}
本文在统一相量、长度和面积定义下，建立了单圆线解析、二维共端场路与有限长度纵向 PEEC 的验证关系。二维内场解的最细网格电阻解析误差为 1.79×10⁻⁶。对端部并联的股线，等流条件可表述为耦合阻抗与外场感应项在差模子空间中的投影为零；640 股直束的复电流结果表明，共同端接与细股划分仍需结合磁耦合评价均流。纵向 PEEC 在已检查长度和积分设置下达到相应敏感性目标，5024 自由度末档电阻与本档等面积 Bessel 参照相差 0.044582\%。

% block-id: 05.conclusions.p2
\hypertarget{05.conclusions.p2}{}
这些结果为后续真实换位结构研究提供了局部场基准、共端约束和有限长度计算的参照。RT0 外场、CPU/GPU 对照及三维诊断资料补充说明了各实现的范围。实际 Litz 结构的性能排序仍需具有端区和横向自由度的三维验证，并结合温升与受控实验完成评价。

% block-id: 05.current.evidence
\hypertarget{05.current.evidence}{}
当前设备的环境迁移重跑了单圆线解析校准与小型 Maxwell 端口及功率案例。这些运行验证迁移后的执行路径，保留原 Maxwell 网格的收敛范围和 A–φ 规范条件失败状态。新核实的半解析与绕组损耗文献为后续工作提供候选方法；其中绕组局部场与均匀化假设，需要与本文孤立导线和共端模型明确比较后，才能转用定量结果。\cite{C01,ZHU2026}
